% SOURCE docs/release/2026-10-29/source/docs/writing-repair/2026-10-29/src/26ee80caee96ff30.tex CHAR 0
\chapter{Four banking translations before formal causal analysis}
\label{exec13:l2-translation}

The four applications first distinguish plausible business questions before
choosing their formal targets. The remainder develops what those translations
must preserve: monetary units and repricing, measurement error and selective
notes, later servicing actions, and the applicant population. Studies of
linguistic measurements, variable construction, model representations and
clinical graph refinement then test the same discipline at more difficult
interfaces. Their different evidence cannot be collapsed into one score for
translation quality.

A request such as ``Should we improve retention?'' is not yet a causal question. The action may be a rate offer, a service contact, a product change or a budgeted policy. The unit may be an account, customer, household or branch. The outcome may be balance retention, funding cost, repayment or relationship value. An assistant should expose these alternatives and ask for the missing decision; a formal engine cannot repair a target chosen incorrectly before the graph is built. The four translations below keep the business motivation while making different estimands visible. Their formal discipline follows the separation of a target trial, its eligible population and its observed emulation in \citet{hernan2020}; the banking choices themselves are proposals, not observations.

\section{Deposit pricing: acquisition, retention or net funding cost}

Suppose a treasury manager considers a 25-basis-point deposit-rate increase. Translation A asks for the effect of offering the rate to every eligible account on thirty-day average balance. Translation B asks for the effect on balances among accounts that historically received similar offers. Translation C asks for net funding cost, including paid interest, incentive expense, repricing of existing balances and any value assigned to retained relationships. Translation D asks for a fixed-budget allocation policy that chooses which accounts receive the offer. These are four populations or outcomes, not four estimators of one effect.

Let $U$ be the account, $P$ the accounts eligible under the written rule, $A$ the offer, and $Y^B(a)=H^{-1}\sum_{d=1}^H B_d(a)$ the average daily balance over $H=30$ days, measured in dollars. The corresponding balance-time outcome is $S(a)=H Y^B(a)$, measured in dollar-days. The all-eligible balance target is $\mathbb E_P[Y^B(1)-Y^B(0)]$. A monetary value target instead specifies a valuation $M(a)$ in dollars and compares $\mathbb E_P[M(1)-M(0)-\{C(1)-C(0)\}]$, where $C(a)$ is monetary cost outside $M(a)$. Costs already inside $M(a)$ must not be subtracted again. If $Y^C(a)$ already contains net funding cost, the objective is to minimize $\mathbb E_P[Y^C(a)]$ and a negative $\mathbb E_P[Y^C(1)-Y^C(0)]$ favors the offer. The policy target additionally requires a feasible policy class and an explicitly costed comparator. A realized post-offer balance may enter the monetary outcome but is not automatically a pre-offer policy weight. A customer identifier does not establish economic ownership when accounts are shared or linked. Section~\ref{exec17:l2-monetary} derives a dimensionally consistent valuation and its alternatives.

For example, let $W$ be the frozen pre-offer balance and segment record, $e(W)=\Pr(A=1\mid W)$, and $m_a(W)=\mathbb E(Y^B\mid A=a,W)$. If consistency, conditional exchangeability, positivity and the declared no-interference boundary hold, the all-eligible target has the identifying expression
\[
 \psi_P=\int\{m_1(w)-m_0(w)\}\,dP(w).
\]
This equality is conditional on the assumptions; a fitted model cannot verify them. If offers are absent for a high-balance stratum, restricting the integral to $0<e(W)<1$ changes its averaging measure. Likewise, a policy $\pi(W)$ under a monetary expense budget $B$ asks for $\mathbb E_P\{Y^C(\pi(W))-Y^C(\pi_0(W))\}$ subject to $\mathbb E_P[C(\pi(W))]\leq B$. The target needs a specified comparator $\pi_0$, a cost convention and an availability clock before any optimization result has a business meaning. The target-versus-estimator distinction is reproduced in the project reference \citep{exec13_formulation_targets}.

A valid clock records when rate eligibility, customer status and balance information were available to the offer system. A statement generated after the decision cannot be used as if it were known before it. Diagnostics include overlap by eligibility stratum, repricing spillovers, concurrent campaigns and whether the offer itself changes future eligibility. A narrower overlap target can be useful, but the report must display the excluded accounts and say that the original all-eligible target remains unresolved.

Consider two equally common eligible strata with conditional effects of $2$ and $8$ balance units. Their all-eligible effect is $5$. If the offer probabilities are $0.1$ and $0.5$, overlap weights $e(1-e)$ are $0.09$ and $0.25$; their normalized average is $(0.09\cdot2+0.25\cdot8)/(0.09+0.25)=109/17\approx6.41$, not $5$. This is exact authored arithmetic, illustrating why a desirable stable-overlap answer cannot be substituted for the requested all-eligible effect. The example asserts no bank offer probabilities.

\section{Arrears servicing: offer, acceptance or dynamic policy}

A collections manager asks whether an early restructuring review should be offered to an account entering arrears. Translation A targets an offer at the first qualifying arrears date. Translation B targets acceptance among customers who take up the review. Translation C targets the effect of a fixed sequence of contacts. Translation D targets a dynamic policy that uses evolving hardship and payment history. Acceptance is downstream of the offer; it is not the same intervention. A static binary ATE cannot silently answer a longitudinal policy question.

Let $X_t$ contain information available at the decision time, $A_t$ the offer, $L_{t+1}$ later hardship information and $Y$ discounted net recovery over 180 days. A total-effect question should not condition on a post-offer descendant $L_{t+1}$ as an ordinary baseline covariate. A policy question needs the action history, availability of each $L_t$, censoring and a comparator policy. If several accounts share a household or collection team, account rows are not automatically independent units. A field entered on day 11 cannot govern a day-10 action even if its event date is day 9.

For a two-decision example let $H_0=X_0$ and $H_1=(X_0,A_0,L_1)$, with a regime $g=(g_0,g_1)$ where $A_0=g_0(H_0)$ and $A_1=g_1(H_1)$. Its target is $V(g)=\mathbb E[Y(g)]$ on the declared first-arrears cohort. Sequential exchangeability, consistency, positivity along the regime and a justified censoring law are needed for a longitudinal identification formula. For the discrete case, write $h=(x,g_0(x),l)$, $p(x)=\Pr(X_0=x)$ and
\begin{align*}
 m_g(h)&=\mathbb E[Y\mid H_1=h,A_1=g_1(h)],\\
 p_g(l\mid x)&=\Pr(L_1=l\mid X_0=x,A_0=g_0(x)).
\end{align*}
Then the longitudinal formula is
\[
 V(g)=\sum_{x,l}m_g(x,g_0(x),l)\,p_g(l\mid x)\,p(x).
\]
The formula illustrates the clock: $L_1$ may guide $A_1$, but cannot retroactively guide $A_0$. It also shows why a coefficient from a regression that conditions on $L_1$ cannot be relabeled the total first-offer effect. This is a project specialization of the longitudinal identification argument in \citet{hernan2020}.

Failure cases are substantive. A model that predicts acceptance can be accurate while answering no offer-effect question. A model that trims accounts with poor overlap and continues to report an all-account effect changes the population. A missing policy-history component is an unsupported target, not a zero effect. The remedy is to return a typed alternative or refusal and request the missing policy, clock or population decision.

\section{Credit underwriting: approval threshold or feasible decision rule}

A credit team may ask whether lowering an approval threshold changes loss. Translation A is the effect of offering credit to applicants between the old and new thresholds. Translation B is the effect on repayment conditional on approval. Translation C is the effect on portfolio loss including volume, utilization and capital cost. Translation D is a policy value under a capacity or risk budget. Approval is an action selected by a threshold, while repayment is observed only under selected applicants; conditioning on approved cases changes the target.

The unit is an application at the decision date. The action must specify amount, term and price, not only a binary approval label. Potential outcomes include repayment, loss, utilization and capital usage over a fixed horizon. A threshold change can alter composition, interference through capacity and future collection intensity. The information contract must distinguish bureau vintage, application-time fields, later delinquency and revised labels. Diagnostics include support around the threshold, manipulation or anticipatory behavior, policy feasibility, missing outcomes and whether the threshold is a rule or a human override.

A regression discontinuity or instrumental design requires its own continuity, exclusion and assignment assumptions; a predictive score does not supply them. If the target is an approval-policy value, the report should retain the policy constraints and loss function rather than output a conditional repayment coefficient. A public or generated threshold mechanism can test translation and clock preservation before requesting bank outcome truth.

To see the selection issue, let $D(z)$ denote approval under threshold policy $z$ and let $Y(1)$ denote repayment under a specified offered loan. The mean $\mathbb E[Y(1)\mid D(z)=1]$ describes the applicants selected by $z$; it is not $\mathbb E[Y(1)]$ for all applicants. Changing $z$ can change that conditioning population even if no individual repayment response changes. Near a score cutoff $c$, the discontinuity $\lim_{s\downarrow c}\mathbb E(Y\mid S=s)-\lim_{s\uparrow c}\mathbb E(Y\mid S=s)$ may have a local causal interpretation only under the design's continuity and manipulation restrictions; it is not the portfolio effect of moving the entire threshold. A threshold policy must count declined applications with their opportunity cost and define the loan terms and missing counterfactual repayment law. The local-design distinction is consistent with the audited discontinuity and target-trial sources \citep{hernan2020}.

\section{Branch campaign: donor counterfactual or randomized rollout}

A branch manager asks whether a campaign increased new accounts. Translation A is a single-branch time-series counterfactual using unaffected donors. Translation B is a panel effect across branches with staggered launch dates. Translation C is a cluster-randomized rollout. Translation D is the value of changing campaign intensity under a constrained budget. The outcomes, donor assumptions and uncertainty objects differ.

For a treated branch $Y_t$, donor series $X_t$ and intervention date $T_0$, a structural time-series method predicts $Y_t(0)$ from pre-period history and controls. It requires controls unaffected by the campaign and a stable relationship after $T_0$. Perfect pre-period fit does not establish either restriction. Shared marketing, customer migration or national pricing can contaminate donors. A randomized rollout addresses a different question with a different cost and design; panel methods require their own assignment and trend assumptions.
The inspected structural-time-series and synthetic-control sources use unaffected predictors and design information in their respective comparisons \citep{exec13_ci_brodersen2015,abadie2010}; neither makes a pre-period fit diagnostic a proof of post-period donor invariance.

An exact construction makes the risk concrete. Let $A_t=1$ on four campaign days, $X_t=x_t^0+15A_t$ and $Y_t=100+2x_t^0+20A_t$. The pre-period relationship is $Y=100+2X$. Continuing it post-period gives a contrast of $-10$ per day, although one structural law has a true campaign effect of $20$ and a contaminated donor. A second law produces the same twelve observed rows and a true effect of $-10$ with an unaffected donor. No estimator can distinguish these laws without a restriction or experiment. With a stipulated USD~1 value per application and USD~20 cost, the four-day net values are +60 and -60. These are authored mechanisms, not bank evidence.

In the first law, the observed post-period contrast is $Y_t-(100+2X_t)=20-2(15)=-10$, while the counterfactual holding the donor's untreated path fixed is $Y_t(0)=100+2x_t^0$ and the true effect is $20$. In the second law, define an unaffected donor $\widetilde x_t^0=X_t$ and an outcome $Y_t=100+2\widetilde x_t^0-10A_t$. Every observed $(A_t,X_t,Y_t)$ row is then identical across the two laws, including the eight pre-period and four post-period rows, yet the effect is $-10$. At USD~1 per application and a single USD~20 campaign cost, $4(20)-20=60$ and $4(-10)-20=-60$. The sign of the investment decision is therefore not identified from these rows alone. The contrasting laws are reproducible algebra, not a claim that either mechanism occurs at a bank.

\section{The formulation record and formal handoff}

The declaration uses the same $Q=(U,P,A,\pi,\pi_0,Y,H,\nu,\ell,I)$ as the preceding formulation chapter: unit, population, feasible actions, policy, comparator policy, outcome, horizon, averaging measure, loss and information clock. A fixed action is a constant policy; $A_0$ in the servicing example denotes the first action, not a second notation for the comparator. It should preserve candidate questions when the business sentence is ambiguous. A schema check verifies fields and types; a semantic check compares them with the brief and dictionary; a formal check applies the supplied graph and method. A passed schema is not causal truth, and a failed backend is not nonidentification.

The order of handoff matters. First, freeze $Q$ and the candidate alternatives with a human owner. Second, attach each proposed graph arrow to a source, temporal ordering and an uncertainty label. Third, ask the formal engine which estimand, if any, follows from the supplied graph and assumptions. Fourth, run estimators only for supported targets and return precision, sensitivity and refusal states. An arrow suggested by retrieval is a hypothesis until its substantive owner accepts it; conversational consistency is not graph accuracy \citep{v12_la_2305_00050,exec13_hallucination}. Formal reasoning benchmarks and narrated synthetic cases help check the translation boundary, but their scores do not validate bank assignment truth \citep{v12_la_2312_04350,v12_la_2607_08093}.

The deterministic baseline is an analyst interview plus fixed arithmetic and graph fixtures. An assisted system can retrieve evidence, draft alternatives, propose a graph and prepare typed calls. It should be scored on target fidelity, clock fidelity, preserved ambiguity, valid abstention, numerical correctness only when the target is identified, and complete failure accounting. All requests belong in the denominator, including refusals and unsupported policies. Public data and generated structural mechanisms are appropriate proposal evidence; authentic bank truth and live provider performance remain later conditional evidence.

The immediate prototype should compare the four translations against the structured baseline using fixed prompts, scenario-level splits, adverse cases and a stopping rule. A later bank study is justified only if the prototype isolates a decision-relevant uncertainty and makes its value of information plausible. No fluent report can substitute for a population decision, assignment record, clock or expert acceptance.

The following studies develop these requirements rather than add another
unqualified promise of automation. They begin with the difference between a
causal sentence and an intervention, then show how text measurement, valuation
and policy definitions change the answer. The final comparisons examine
whether proposed language and graph components preserve those distinctions.


% SOURCE docs/release/2026-10-29/source/docs/writing-repair/2026-10-29/src/4396cbc18328130c.tex CHAR 0
\section{Contextual causal language and the intervention a bank will buy}
\label{exec23:context-scope}

A causal sentence in a service note is useful evidence about what was said,
but the intervention a bank will buy requires an additional scientific
argument. Zhao, Yan and Hu's \emph{Contextual Causality with Large Language
Models: A Survey}, inspected in its 18 September 2026 arXiv version, separates
semantic, intervention and counterfactual tasks within textual context. Its
page 3 explicitly distinguishes intervening in a textual description from
manipulating the underlying data-generating process. The paper's semantic
``identification'' of causal relations should therefore be read as a language
task, rather than statistical identification of an effect. The authors state
that their text-only taxonomy organizes causal concepts and does not extend
causal theory.~\citep{exec23_contextual_causality_survey}

For banking translation, the useful outputs are candidate entities, proposed
temporal relations, alternative explanations and the passages supporting each
proposal. An analyst may translate ``customers stayed after we called'' into
an action, outcome and eligible population. The same note may omit that
relationship managers called customers they already expected to retain. A
language model can make both possibilities explicit. A fluent restatement,
or a correct answer to a question about the note, cannot determine which
mechanism generated the observations.

\subsection{The same observed relationship can reverse the investment decision}
\label{exec23:context-twins}

An exact authored pair of causal laws isolates the missing information. Let
$U\sim\operatorname{Bernoulli}(1/2)$ be an unobserved account state, and let the
natural binary contact rule be $A=U$. Let $Y(a)$ be gross monetary response
under assigned contact $a$. Consider
\begin{equation}
 \mathcal M_1:\ Y(a)=100a,
 \qquad
 \mathcal M_2:\ Y(a)=100U.
 \label{exec23:context-laws}
\end{equation}
Both laws generate exactly
\begin{equation}
 P(A=0,Y=0)=\tfrac12,
 \qquad P(A=1,Y=100)=\tfrac12.
 \label{exec23:context-observed}
\end{equation}
In $\mathcal M_1$, contact causes the gross response. In $\mathcal M_2$ the
latent state causes both natural contact and response, and contact has no
effect. Under an intervention the distribution of $U$ remains unchanged.
Hence the average gross effect is USD 100 in the first law and zero in the
second.~\citep{exec23_context_reference}

At USD 10 per contact, changing from no contact to contacting every eligible
account has incremental net value
\begin{equation}
 E\{Y(1)-Y(0)\}-10=
 \begin{cases}
 90,&\mathcal M_1,\\
 -10,&\mathcal M_2.
 \end{cases}
 \label{exec23:context-net}
\end{equation}
All values and mechanisms here are stipulated; they are not bank observations
or outputs of an LLM. The target compares two fixed policies, all-contact and
no-contact, rather than comparing either policy with the natural contact rule.
The witness proves that the displayed observed law alone cannot identify this
decision. It does not prove that every useful causal question is unidentified.

An assistant should preserve both graph possibilities when the supplied
evidence cannot distinguish them. A new randomized contact assignment with
adequate support would distinguish these two gross-response laws under its
design assumptions. Observing the latent state alone does not create the
missing within-state contact variation because $A=U$ under the natural rule.
Alternatively, a justified exclusion of the latent common-cause explanation
would change the model class, but that restriction needs evidence. Repeatedly
asking the same model to rewrite the same observations does not add either
source of information. The proposed business-to-graph flow should report the
unresolved distinction and the experiment or restriction that could resolve it.

\subsection{A survey can guide the evaluation design without qualifying every cited method}
\label{exec23:context-evaluation}

The survey discusses hallucinated or neglected relations, complex contexts,
multiple-choice shortcuts and limitations of external knowledge and agent
coordination. These are the survey authors' synthesis of their cited work;
those originals require separate inspection before their empirical claims
can support a package decision. Table 3 draws many representative benchmarks
from news, Wikipedia, stories, human questions and agent traces, and includes
FinCausal 2020. The presence of financial text is not evidence of contact-policy
identification or monetary utility in banking. Pages 8--9 call for extrinsic
decision quality, finer evaluation and longer agent horizons. This is a
research agenda, not an authenticated exhaustive absence result.
\citep{exec23_contextual_causality_survey}

A public or generated prototype can evaluate the resulting proposal before
bank access. Supply notes with controlled omissions, negation, reversed
chronology and conflicting explanations; separately score relation extraction,
clock preservation, graph alternatives, unsupported restrictions and appropriate
refusal. Pair identical observations from the two laws above and check that an
assistant does not invent information to force a unique effect. When an
assignment design or restriction is supplied, evaluate the changed answer and
its stated dependence on that premise. These proposed tests assess different
tasks from factual question answering. Their known-law monetary scoring is
conditional on the authored mechanisms and cannot establish actual bank value.


% SOURCE docs/release/2026-10-29/source/docs/writing-repair/2026-10-29/src/9434581a1e2ca284.tex CHAR 0
\hypertarget{text-measurement-must-preserve-the-causal-question}{%
\subsection{Text measurement must preserve the causal
question}\label{text-measurement-must-preserve-the-causal-question}}

\label{exec13:census-measurement}

A servicing analyst may ask an assistant to extract financial distress
from call notes and then adjust for the extracted score when comparing
repayment plans. This can be useful work. It also combines three
different claims: that the score measures distress, that the note was
available before the plan decision, and that conditioning on the score
removes the relevant confounding. An accurate summary of the note
establishes none of the latter two claims. A later note may describe a
reaction to the plan; a note written only after a complaint may select
on both treatment and outcome. The proposed formulation interface must
attach a time and a measurement interpretation to each extracted
variable before adding it to an adjustment set.

Keith, Jensen and O'Connor review text used to measure otherwise
unobserved confounders. Their discussion explicitly distinguishes
confounders, colliders and mediators and asks whether one document can
encode several of these at once. It also identifies error in measured
confounders as an unresolved problem. A richer language representation
can improve measurement, but the review provides no guarantee that a
representation eliminates confounding. The appropriate comparison
therefore includes simple dictionaries, supervised extractors,
human-coded measurements and a language model, with the same permitted
document dates and the same effect target.
\cite{exec13_census_text_review_arxiv}

\hypertarget{a-reproducible-example-of-residual-confounding}{%
\subsubsection{A reproducible example of residual
confounding}\label{a-reproducible-example-of-residual-confounding}}

Consider the following project-generated population model. The
intervention intensity \(A\) and later loss \(Y\) are standardized
continuous quantities; these are illustrative units, not calibrated bank
measurements. True predecision distress is \(U\), while \(W\) is an
imperfect score extracted from a predecision note:

\[
 A=U+\varepsilon_A,\qquad
 W=U+\varepsilon_W,\qquad
 Y=\tau A+U+\varepsilon_Y.
\]

All four disturbances are mutually independent and centered Gaussian.
The variances of \(U,\varepsilon_A,\varepsilon_Y\) are one, while
\(\operatorname{Var}(\varepsilon_W)=s\geq0\). The true effect of
increasing \(A\) by one is \(\tau\). Conditional on \(W\), the residual
variance of \(U\) is \(v=s/(1+s)\). Residualizing both \(A\) and \(Y\)
on \(W\) therefore gives the population regression coefficient

\[
 \beta_{A\mid W}
 =\tau+\frac{v}{1+v}
 =\tau+\frac{s}{1+2s}.
\]

The second term persists with an infinite sample. It is not estimation
noise that a larger language model, an accurate propensity fit to \(W\),
or a narrower regression interval automatically removes. With the
stipulated \(\tau=-1/5\), exact covariance calculations give the
following values.

\begin{longtable}[]{@{}lll@{}}
\toprule
Text-score noise variance \(s\) & True effect & Adjusted population
coefficient\tabularnewline
\midrule
\endhead
\(0\) & \(-1/5\) & \(-1/5\)\tabularnewline
\(1/4\) & \(-1/5\) & \(-1/30\)\tabularnewline
\(1\) & \(-1/5\) & \(2/15\)\tabularnewline
\(4\) & \(-1/5\) & \(11/45\)\tabularnewline
\bottomrule
\end{longtable}

The score with \(s=1\) leaves enough confounding to reverse the sign.
This example is independently reproduced by solving the population
covariance normal equations and by residual variance decomposition. It
demonstrates a specific failure mechanism under a stated law, not a
claim that all text adjustment fails. A successful extraction study must
measure the relevant error and evaluate how it affects the causal
target. If a proxy-correction method is proposed, its additional
measurement assumptions and available validation information must be
stated rather than inferred from predictive accuracy.

\hypertarget{freeze-the-measurement-rule-before-evaluating-its-effects}{%
\subsubsection{Freeze the measurement rule before evaluating its
effects}\label{freeze-the-measurement-rule-before-evaluating-its-effects}}

Egami and colleagues make the measurement rule itself explicit. Let
\(g\) map a document into a category or vector. For a textual outcome,
\(E[g(Y(1))-g(Y(0))]\) is a target defined by that particular rule.
Searching over rules using the same treatment assignments and documents
later used for effect estimation can induce dependence and select
apparent effects. Their proposed separation of codebook discovery from
effect estimation addresses that problem, while their text-as-treatment
discussion imposes a further sufficiency restriction on the
representation. \cite{exec13_census_texts}

For a language-based measurement pipeline, the rule includes the model
version, prompt, retrieval corpus, preprocessing, extraction schema and
response aggregation. Freezing only a classifier filename while changing
the prompt after reading the effect estimates does not freeze \(g\). The
prototype should develop the rule on a designated development set,
record its identity, and evaluate it on held-out units. If several calls
concern one customer or conversation, those calls should remain together
in the split. This is a proposed implementation of the separation
principle, not a claim that sample splitting alone ensures valid
inference for dependent histories.

The held-out review asks whether the labels retain their intended
meaning, whether extraction errors differ by treatment and outcome, and
whether missing documents have a defensible observation model. It must
preserve negative findings. Rewriting the measurement rule to rescue a
disappointing effect requires a new evaluation set or an explicitly
exploratory analysis. Neither a fixed rule nor a held-out set supplies
unmeasured assignment information or repairs a post-treatment adjustment
error.

\hypertarget{formulate-competing-explanations-before-selecting-a-diagram}{%
\subsubsection{Formulate competing explanations before selecting a
diagram}\label{formulate-competing-explanations-before-selecting-a-diagram}}

Geraldo Bastías proposes building diagrams from the assumptions actually
used in a study, preserving several models where the assumptions do not
determine one. His workflow checks identification within each model,
derives observable implications, and distinguishes missing relationships
from missing variables. This is a methodological contribution to diagram
review, not an implemented LLM discovery algorithm. It supports an
assistant that makes assumptions inspectable rather than treating a
familiar design label as sufficient. \cite{exec13_census_credible}

For the servicing example, a short initial brief can leave open whether
distress caused the plan assignment, whether a later complaint caused a
note to be recorded, and whether the proposed outcome is arrears, loss
or customer welfare. The assistant should construct separate compatible
diagrams and show how each changes adjustment or identification. It can
check timing contradictions and derive conditional-independence
implications. Passing those checks does not establish that an omitted
common cause is absent. Where the diagrams imply different effects, the
useful output is the unresolved choice and the evidence that could
distinguish it.

These requirements permit immediate public and generated evaluation. A
public document corpus can supply realistic language, while a declared
generator supplies assignment, document timing, measurement error and
outcome truth. Simple extraction and fixed numerical workflows provide
baselines. The evaluation must vary mechanisms that are difficult for
the chosen extractor, including selective note availability and notes
written after treatment. Keith and colleagues specifically caution that
authors can choose semi-synthetic confounders their own method can
easily recover. Freezing adverse mechanisms and separating their
construction from final method selection addresses that concern more
directly than adding familiar domain vocabulary to a favorable
simulation. \cite{exec13_census_text_review_arxiv}

\label{exec13:census-measurement-end}


% SOURCE docs/release/2026-10-29/source/docs/writing-repair/2026-10-29/src/bcd8d1d450eacae8.tex CHAR 0
\section{From a balance response to a monetary decision}
\label{exec17:l2-monetary}

The deposit example exposes a difficulty that arises before any estimator is
chosen. A balance is a stock, its accumulation over time is an exposure, and
the financial benefit of that exposure depends on a valuation convention.
These quantities cannot be interchanged because they all contain a currency
unit. The assistant must obtain a complete statement of the business target:
does treasury want additional deposits, lower expense for a fixed funding
requirement, or a relationship-value objective that includes services beyond
funding? Each question can be legitimate. Only the second is answered by the
funding-cost construction developed here. The construction and all numerical
values in this section are authored, exact illustrations. They supply a
reference for testing a proposed translator, not observed bank economics
\citep{exec17_l2_monetary_exact}.

Let $B_d(a)$ be end-of-day eligible balance under offer $a$ on day $d$, where
$a=0$ is the stated existing offer and $a=1$ is the proposed offer. Fix a
horizon of $H$ days and a day-count denominator $D$. Define
\begin{equation}
 \overline B(a)=\frac1H\sum_{d=1}^H B_d(a),\qquad
 S(a)=\sum_{d=1}^H B_d(a)=H\overline B(a).
 \label{exec17:l2-stock}
\end{equation}
The units are dollars and dollar-days respectively. With a constant annual
valuation rate $q$, the monetary amount assigned to the balance-time is
$qS(a)/D$ dollars. The rate, denominator, daily measurement convention and
horizon are part of this definition. For example, substituting a last-day
balance for an average balance changes the outcome even if the two amounts
happen to agree in a stable illustrative account. A genuine proposal needs
the product's actual accrual convention; the example fixes $D=365$ without
claiming that this convention applies to every deposit contract.

Suppose, for the moment, that treasury has the same exogenous daily funding
requirement $F_d$ under either offer and can cover the residual $F_d-B_d(a)$
at a supplied rate $q_d$. This construction assumes $0\leq B_d(a)\leq F_d$; excess deposits would require a separately specified investment or disposal rule. Let $r_d(a)$ be the deposit rate, and let $K(a)$ be
the monetary expense outside interest: for example a stipulated communication
charge. Ignoring discounting in this short-horizon construction, define
\begin{align}
 N(a)&=\frac1D\sum_{d=1}^H
 \{r_d(a)B_d(a)+q_d[F_d-B_d(a)]\}+K(a),
 \label{exec17:l2-cost}\\
 V(a)&=\frac1D\sum_{d=1}^H[q_d-r_d(a)]B_d(a)-K(a).
 \label{exec17:l2-value}
\end{align}
Here $N$ is net funding cost and $V$ is the saving relative to funding the
entire requirement at $q_d$. Their sum is the common reference cost
$D^{-1}\sum_d q_dF_d$. Therefore the offer's incremental value is exactly
$\Delta V=\mathbb E[V(1)-V(0)]=-\mathbb E[N(1)-N(0)]$.
Maximizing $V$ and minimizing $N$ agree. Neither convention licenses
maximizing the increase in funding cost. A separate expense must be netted
as $K(1)-K(0)$, not automatically as $K(1)$: the comparator itself may incur
expense, and cancelling a campaign may save it.

This identity depends on the reference funding requirement and valuation
path being common. If the offer changes lending volumes, the institution's
marginal funding rate, or the capacity to raise substitute funds, those
objects acquire potential-outcome indices and the common-reference term no
longer cancels. An assistant cannot preserve the simple identity while
quietly adding those effects to its narrative. It must expand the target,
state the additional intervention and equilibrium assumptions, or present
the fixed-requirement result as a narrower answer. Likewise, relationship
benefits should be represented by a separately specified monetary outcome,
not attached to a balance effect through an unexplained positive adjective.

\subsection{Repricing existing balances can reverse the conclusion}

Under constant rates, write $r_1=r_0+\delta r$ and
$\Delta S=S(1)-S(0)$. Subtracting the two value definitions gives the
pathwise identity
\begin{equation}
 V(1)-V(0)=\frac{(q-r_0)\Delta S-\delta r S(1)}D
             -[K(1)-K(0)].
 \label{exec17:l2-decomposition}
\end{equation}
The first term values additional balance-time at the old spread. The second
charges the higher deposit rate to all balances receiving it, including
balances that would have remained without the offer. It is not restricted
to newly attracted deposits. An equivalent decomposition uses the new
spread on $\Delta S$ and the rate increase on $S(0)$. Both produce the same
cash difference; adding both repricing terms would count part of the cost
twice. These identities hold for each pair of potential balance paths, so
taking expectations requires no independence between the two paths.

Consider an exact thirty-day example. The comparator holds USD~100,000 every
day at a 2\% annual rate; the offered rate is 2.25\%, and its balance is
USD~110,000 every day. The stipulated substitute funding rate is 5\%.
Operating expenses are USD~4 under the comparator and USD~8 under the offer.
Thus $\Delta S=300{,}000$ dollar-days and the added average balance is
USD~10,000. The additional old-spread value is $1800/73$ dollars and the
repricing charge is $1650/73$ dollars. Their difference is only $150/73$
dollars. Subtracting the incremental USD~4 expense gives
\begin{equation}
 \Delta V=\frac{1800-1650}{73}-4=-\frac{142}{73}
          \simeq -\mathrm{USD}\,1.95.
 \label{exec17:l2-numeric}
\end{equation}
The offer attracts balances while increasing net funding cost. The reference
also constructs the complete daily ledger for a stipulated USD~150,000 funding
requirement and reproduces the opposite-signed cost difference by summation.
That second calculation checks the accounting identity instead of merely
repeating the printed formula. Treating the USD~10,000 stock increase as a
USD~10,000 benefit, forgetting existing-balance repricing, or subtracting
USD~8 instead of the incremental USD~4 each answers a different question.

The all-eligible mean effect remains a population target. A deterministic
balance path here makes the arithmetic transparent; a causal analysis must
identify or bound the distribution of paths under both offers. Identifying
the mean balance contrast alone is sufficient for this linear, fixed-rate
valuation only when the other required expectations and costs are supplied.
It is not a claim that a regression coefficient identifies the complete
financial response of a deposit portfolio.

\subsection{Marginal funding cost and an explicit valuation boundary}

Another reasonable business question is the annualized incremental expense
per unit of extra balance-time. For a positive expected balance-time change,
under the same constant-rate convention define
\begin{equation}
 r_{\mathrm{inc}}=
 \frac{r_1\mathbb E[S(1)]-r_0\mathbb E[S(0)]
       +D\mathbb E[K(1)-K(0)]}
      {\mathbb E[S(1)-S(0)]}.
 \label{exec17:l2-marginal}
\end{equation}
This is a ratio of portfolio expectations, not the mean of customer-level
ratios. The latter can be undefined for accounts with no balance response,
give large influence to tiny denominators, and change the aggregation
question. The incremental cost rate also differs from the offered deposit
rate. In the example it is $1571/30000$, or about 5.237\%, despite the
2.25\% offered rate. It includes paying more on existing balances and
spreading the incremental operating expense over the new balance-time.
With positive expected $\Delta S$, $\Delta V$ is positive precisely when
$q>r_{\mathrm{inc}}$. If that denominator is zero or negative, the displayed
ratio is unavailable or loses this comparison interpretation; the original
monetary difference remains the appropriate object to inspect.

This comparison is useful even when a bank has not yet supplied a valuation
rate. Instead of inventing a price, the proposal can report the break-even
rate or a sensitivity interval. If the same example permits $q$ between
4\% and 6\%, its value ranges from $-742/73$ to $458/73$ dollars, roughly
minus USD~10.16 to plus USD~6.27. There is no uniformly preferred offer over
that stipulated set. An owner who accepts the funding model can resolve
the economic choice by supplying the relevant rate and cost convention.
This is a missing business input, distinct from an unidentified balance
response. Uncertainty about both should be kept visible rather than folded
into one unexplained error bar.

For a set of admissible causal laws $\mathcal P$ and valuation specifications
$\mathcal Q$, the corresponding proposal can report
\begin{equation}
 \mathcal V=\{\mathbb E_P[V_q(1)-V_q(0)]:(P,q)\in\mathcal A\},
 \qquad \mathcal A\subseteq\mathcal P\times\mathcal Q.
 \label{exec17:l2-joint}
\end{equation}
The joint admissible set $\mathcal A$ matters. If liquidity conditions affect
both the response and the marginal rate, independently combining marginal
endpoints may create impossible scenarios or omit admissible dependence.
The interval above deliberately varies only a deterministic rate while
holding the response law fixed. It is an economic sensitivity calculation,
not a confidence interval, a Bayesian posterior, or a general sharp bound.
The formulation should state which of those objects the user is requesting
before selecting a downstream numerical tool.

\subsection{Timing and nonlinear valuation require more than a mean}

Balance-time suffices only for a constant linear valuation. With a supplied
time-varying path $q_d$, the gross monetary effect is
$D^{-1}\sum_d q_d\mathbb E[B_d(1)-B_d(0)]$. As an exact two-day example,
compare additional balances of $(200,0)$ with $(0,200)$ dollars. Both have
average additional balance USD~100 and exposure 200 dollar-days. At rates
$(0.10,0.01)$ their gross values are $20/365$ and $2/365$ dollars. A translator
that passes only the mean balance to the valuation engine loses the
information needed to distinguish these cases. Discount factors, when
required, must similarly attach to the appropriate cash-flow dates; they
cannot be appended after aggregation if the timing information was dropped.

A nonlinear valuation needs still more information. Suppose a hypothetical
gross value rule is $m(B)=0.01\min(B,100)$ dollars and a comparator has zero
gross value. A balance fixed at USD~100 has expected gross value USD~1.
A balance of zero or USD~200 with equal probabilities has the same mean
balance but expected gross value USD~0.50. With incremental expense USD~0.75,
their net values are USD~0.25 and minus USD~0.25. The mean balance response
does not identify which decision is favorable. This example is deliberately
stylized; its purpose is to distinguish $\mathbb E[m(B)]$ from
$m(\mathbb E[B])$, not to suggest a bank's valuation schedule.

Where costs depend on utilization, default, or other post-offer quantities,
those quantities belong in the monetary potential outcome. Their joint
distribution can matter. Using the realized balance to calculate interest
is legitimate outcome construction; using that same post-offer balance as
a pre-offer targeting feature is a clock violation. Similarly, conditioning
on it to estimate an ordinary total effect generally changes the question.
The declaration must identify the role a variable plays, rather than
labelling a field universally permitted or forbidden. The target-trial and
time-varying treatment distinctions already developed from
\citet{hernan2020} provide the inherited framework; the cash-flow mappings
here are project-specific definitions.

\subsection{Consolidation and feasibility are part of the intervention}

The business unit determines what counts as additional funding. If an offer
moves USD~10,000 from one account to another account at the same bank, the
receiving account's effect is positive while the consolidated balance change
is zero. Repricing the receiving account can nevertheless increase expense.
For mutually exclusive account balances under a fixed consolidation map,
the bank-level outcome sums the account paths before applying the bank-level
valuation. Shared accounts must not be counted once for each linked owner.
When the intervention changes ownership, currency composition or legal
availability, a fixed consolidation map itself becomes an assumption to
reconsider. An identifier matching exercise cannot choose the economic
boundary on the owner's behalf.

This also changes treatment definition. A rate offered to one account may
alter another account's balance or a household's other offers. The potential
outcome may then depend on an allocation vector, not just an account's own
binary action. A feasible policy must respect the shared budget and the
permitted unit of allocation. Account-level exchangeability and a
no-interference assumption are substantive restrictions on that target,
not conveniences a language model can infer from a rectangular dataset.
If these restrictions are not defensible, the proposal can retain a cluster
or portfolio question with its additional requirements.

A budget needs an equally clear comparator. An absolute expense constraint
$\mathbb E[K(\pi)]\leq B$ differs from an incremental constraint
$\mathbb E[K(\pi)-K(\pi_0)]\leq B$. A capacity constraint on the number of
contacts is a third object and may not be expressible as either monetary
constraint. The policy $\pi$ can use only the declared predecision history;
the comparator $\pi_0$ must satisfy its own declared feasibility conditions.
The proposed assistant should present these alternatives in the business
brief and carry the accepted choice into the formal query. Optimization
before that choice can produce a precise answer to an unstated problem.

\section{Measurement and selective notes can change the target}
\label{exec17:l2-measurement}

The cash-flow argument generalizes to language-derived outcomes. A rule
$g$ that converts a servicing note into a distress indicator defines an
observed measurement, not automatically the underlying distress variable.
The inspected discussions by Keith and colleagues and Egami and colleagues
separate measurement from causal identification and distinguish a learned
text representation from the scientific target
\citep{exec13_census_text_review_arxiv,exec13_census_texts}. The proposed
translator must say whether the owner wants the effect on recorded distress,
on a fixed coding rule's output, or on latent distress. Those outcomes can
differ even when the extraction system executes exactly as specified.

Let a binary underlying outcome have intervention-specific mean
$\mu_a=\Pr(Y(a)=1)$. For a fixed binary measurement rule, define its
sensitivity $s_a=\Pr(Y^*=1\mid Y=1,\operatorname{do}(A=a))$ and specificity
$c_a=\Pr(Y^*=0\mid Y=0,\operatorname{do}(A=a))$. By the law of total
probability, when these conditional rates are defined,
\begin{equation}
 \mu_a^*=s_a\mu_a+(1-c_a)(1-\mu_a)
          =(s_a+c_a-1)\mu_a+1-c_a.
 \label{exec17:l2-error}
\end{equation}
This is an exact project derivation, not an imported guarantee for an LLM.
If sensitivity and specificity are common across interventions, the measured
contrast equals $(s+c-1)(\mu_1-\mu_0)$. Recovering the underlying contrast
by division requires those error rates to be known, transportable and
$s+c-1\ne0$. If the sum is one, the measurement contains no marginal
information about the binary outcome through this channel. If it is less
than one, its direction is reversed. Assuming ordinary attenuation without
checking the channel is therefore unjustified.

For example, set true outcome risks to $2/5$ and $1/2$, a true effect of
$1/10$. With common sensitivity $4/5$ and specificity $9/10$, the measured
effect is $7/100$. Now retain those control error rates but use treated
sensitivity $1/2$ and specificity $19/20$. The measured effect becomes
$-21/200$, reversing the sign. All probabilities remain valid. The change
could represent different document content or extraction errors after the
offer; the example makes no claim that this occurs in a bank. Its implication
is that pooled classification accuracy cannot determine the causal error.
A validation design needs outcome truth and enough treatment-specific
information to support the proposed correction, including the population
and document mechanism to which the rates are transported.

Missing notes introduce a separate problem even when every available note
is coded perfectly. Let $A$ and predecision distress $U$ be independent
Bernoulli variables with probability $1/2$, and define $Y(a)=U$ for either
action. The intervention has zero effect. Let $R$ indicate a note available
for analysis, with probability $3/4$ when $A=U$ and $1/4$ otherwise.
Every action--distress combination has positive note availability, so this
example does not depend on a cell with zero support. Exact enumeration gives
\begin{equation}
 \Pr(Y=1\mid A=1,R=1)=\frac34,\qquad
 \Pr(Y=1\mid A=0,R=1)=\frac14.
 \label{exec17:l2-selection}
\end{equation}
The observed-note contrast is $1/2$ despite random assignment and a zero
effect in the full population. Selecting notes conditions on a common effect
of action and distress. A correct event timestamp, perfect extraction and
randomized treatment do not remove that selection mechanism.

An oracle that knows $U$ and the note probabilities can reweight each
selected cell by the inverse of its selection probability and recover
mean distress $1/2$ in each arm. This verifies the construction; it is not
a proposed correction when distress is precisely the unobserved quantity
the notes were meant to measure. The useful handoff identifies what is
missing: a justified observation model, independent validation measurements,
or a narrower target that the owner actually accepts. It must not quietly
replace the eligible population by customers with convenient documentation.

These examples distinguish three repairs. Correcting a unit or comparator
is a specification repair. Estimating a measurement channel needs suitable
validation information. Correcting selective availability needs a defensible
observation law and support. A language model can locate definitions,
propose alternatives and flag contradictions in each case, but those actions
do not provide the missing scientific information. Multiple defensible
diagrams and their different identifying consequences should survive the
formulation stage, consistent with the inspected argument for explicit
assumption-based modelling in \citet{exec13_census_credible}.

\subsection{A proposal evaluation that isolates the translation error}

A useful public and generated evaluation can freeze a set of business
briefs that share language while differing in one decision-relevant fact.
For funding, pair a balance target with a net-cost target; then vary whether
the comparator incurs expense, whether repricing covers old balances, and
whether valuation is time-varying or nonlinear. For measurement, vary
common versus differential error and complete versus selective notes.
The exact mechanisms above supply the expected mathematical answer, while
public dictionaries can supply realistic field names and documented units.
Public text alone supplies neither assignment truth nor a bank's marginal
funding rate. An evaluation should preserve that distinction explicitly.

The reference response should contain the target, the necessary outcome
components, the sign convention, the information that is absent, and the
reason an alternative interpretation would change the answer. A numerical
score is appropriate only when the brief supplies enough information to
determine it. For an unspecified $q$, reporting the break-even rate and the
unresolved valuation choice can be correct; inventing a rate to produce a
confident investment recommendation is an error. For selective notes with
unknown selection probabilities, reporting the selected contrast as the
full-population effect is an error even when the subtraction is exact.

The comparison should use the same briefs, dictionaries, permissible source
material and answer schema for a structured analyst baseline and an assisted
workflow. Separate the underlying mechanisms used for development from those
used for final evaluation, and keep ambiguous, unsupported and failed requests
in the denominator. This section provides exact authored reference families,
not a completed provider comparison. A useful prototype result would identify
which translation failures become less frequent and at what review cost.
That result can support a later decision to purchase authentic banking
validation; it cannot supply such validation in advance.
\label{exec17:l2-monetary-end}


% SOURCE docs/release/2026-10-29/source/docs/writing-repair/2026-10-29/src/9a18115db138948f.tex CHAR 0
\section{Which servicing policy does the business sentence describe?}
\label{exec17:l2-policy}

A request to improve early servicing can leave the later action rule
unspecified. That omission is more consequential than a missing estimator
option. An initial review may change a later condition, and a subsequent
contact may respond to that condition. Comparing first-stage actions while
holding the later action fixed therefore differs from comparing complete
operating policies. The longitudinal formula developed earlier makes this
distinction precise; the following finite construction shows the numerical
consequence. It specializes the inherited sequential-intervention framework
of \citet{hernan2020} to an authored example, with no empirical claim about
how an actual servicing review affects hardship.

Let $A_0$ denote an initial action, $L_1$ a later adverse-state indicator,
and $A_1$ a follow-up action taken after observing $L_1$. Stipulate
$\Pr(L_1=1\mid\operatorname{do}(A_0=0))=1/4$ and
$\Pr(L_1=1\mid\operatorname{do}(A_0=1))=3/4$. The initial action has a
direct benefit but increases the chance of the costly later state in this
deliberately adverse model. Define net recovery, in hypothetical dollars, by
\begin{equation}
 Y(a,l,b)=100+20a-40l+30bl-10b.
 \label{exec17:l2-sequential-reward}
\end{equation}
The last term is the follow-up cost and is already included in the outcome.
There is no additional initial cost in this construction. The structural
model uses an independent uniform disturbance to generate $L_1$ at the
stated probabilities, and the outcome is then deterministic given the
actions and state. These declarations supply the causal law for the example;
observed historical conditional frequencies would not by themselves supply
the same law.

For a fixed initial action $a$, let $g^a$ take the follow-up action precisely
when $L_1=1$. In that state, follow-up adds USD~20 net; when $L_1=0$ it
subtracts USD~10. The adaptive rule is consequently optimal within this
specified two-action second-stage class when the state is fully observed
and there is no shared capacity constraint. Its value is
\begin{equation}
 V(g^a)=\sum_{l=0}^1Y(a,l,l)\Pr(L_1=l\mid\operatorname{do}(A_0=a)),
 \qquad V(g^0)=95,\quad V(g^1)=105.
 \label{exec17:l2-sequential-value}
\end{equation}
The complete-policy contrast is USD~10. If follow-up is prohibited under
both initial actions, both values are USD~90 and the first-action contrast
is zero. If the state is held fixed and the same follow-up action is used,
the direct contrast in the reward expression is USD~20. Finally, comparing
$g^1$ with never taking either action yields USD~15. All four calculations
are correct for their own interventions. A translator that reports one
without stating which subsequent policy is held fixed has not resolved the
business question. Subtracting the follow-up expense again would introduce
a separate accounting error.

The value comparison can be independently checked under a fully supported
generated logger. Let the logger randomize each action with probability
$1/2$, using an independent randomization at the second stage after observing
the state. Its eight possible histories have probabilities
$\tfrac14\Pr(L_1=l\mid\operatorname{do}(A_0=a))$. For the target $g^a$,
retain histories with $A_0=a$ and $A_1=L_1$ and multiply their logger
probabilities by four. Summing the weighted rewards reproduces 95 and 105.
This finite enumeration and the direct policy calculation are implemented
separately in the exact reference \citep{exec17_l2_policy_exact}. The agreement
checks this particular known law. It does not show that a real logger is
known, that hidden hardship is observed, or that weights fitted to dependent
customer histories have calibrated uncertainty.

The availability clock gives the rule an operational interpretation. If
$L_1$ is available before the second decision, $A_1=L_1$ is admissible.
If the note describing it arrives after that decision, the same expression
uses unavailable information. A fixed always-contact rule has values
USD~87.50 and USD~102.50 for initial actions zero and one; a fixed never-contact
rule has values USD~90 and USD~90. The best fixed rule is never-contact for initial action zero and
always-contact for initial action one under this stipulated law. These rules
illustrate a value difference caused by changing the permitted information,
not by switching estimators. A forecast available earlier is another
candidate input with its own joint law; it cannot be substituted for the
later realized state without a new calculation.

The formulation should accordingly distinguish three requests: evaluation
of a supplied policy, selection from a declared feasible policy class, and
an unrestricted optimal-policy claim. The example resolves only a tiny
known-law class. Adding hidden state, nonbinary actions, capacity constraints
or an unknown behavior policy creates additional research and implementation
requirements. The assistant can preserve these requirements in a formal
query and explain why the static effect is insufficient. It cannot convert
the finite optimum into a guarantee for general private-state servicing.

\section{A threshold change is an applicant-population decision}
\label{exec17:l2-threshold}

Credit approval introduces a different translation failure: observing a good
outcome among accepted applicants does not establish the value of expanding
acceptance. Let $D_z(S)$ be a deterministic approval rule at score $S$ under
policy $z$. Define $R(1)$ as net monetary contribution if a specified loan
is offered and $R(0)$ as the contribution under the specified alternative,
such as declining this product while retaining other services. On a common
applicant population, the policy value is
\begin{equation}
 V(z)=\mathbb E[D_z(S)R(1)+\{1-D_z(S)\}R(0)].
 \label{exec17:l2-credit-value}
\end{equation}
This definition needs the loan terms, horizon, default and recovery treatment,
capital charges and the decline alternative fixed. Repayment probability
alone does not supply $R$. Nor is the mean contribution among approved
applications equal to $V(z)$: its denominator changes with the threshold.

Suppose the new rule expands approval monotonically and changes no loan
terms for already approved applicants. Under a common potential-outcome law
without capacity spillovers, subtracting the policy values gives
\begin{equation}
 V(1)-V(0)=\mathbb E[\{D_1(S)-D_0(S)\}\{R(1)-R(0)\}].
 \label{exec17:l2-credit-switchers}
\end{equation}
Only the newly approved band contributes. If repricing or other changes also
affect incumbent approvals, this simplification fails and their contribution
must remain in the target. If actions interact through a capital or staffing
constraint, the individual response law also needs revision. The equation
is an accounting consequence of the specified potential outcomes, not an
identification result from historical approvals.

For an exact counterexample, let the old rule approve one fifth of applicants,
each with USD~100 expected net contribution, and let the new rule admit an
additional one tenth. Set the decline contribution to zero. In one admissible
law the new band would contribute USD~50 each if offered the loan; in another
it would contribute minus USD~50. The historical data under the old rule can
be identical because this band received no loans. Its loan outcomes remain
unobserved. The old policy value is USD~20 per applicant in either law, while
the expanded policy has values USD~25 and USD~15. The incremental values
are therefore USD~5 and minus USD~5. Exact enumeration verifies the result
\citep{exec17_l2_policy_exact}.

Meanwhile, mean contribution per approved applicant falls from USD~100 to
USD~$250/3$ in the favorable law, and to USD~50 in the adverse law. Thus even
an observed reduction in the approved-group average does not alone show that
expansion reduces total applicant-population value. It can reflect adding a
positive but less profitable group. The proposal should carry the common
applicant denominator, the new approval mass and the net contribution
convention through the analysis. A question about whether loan quality falls
should retain portfolio value and conditional credit quality as
different questions until the owner chooses the intended objective.

Suppose only the scientific range
$-200\leq R(1)-R(0)\leq100$ is supplied for the new band. Its known mass
of one tenth then yields the incremental value range $[-20,10]$ dollars
per applicant. Both endpoints are attained by assigning the corresponding
constant contribution throughout the unobserved band, without changing the
old-policy observations. This is a sharp range within the stipulated class,
not a sampling interval or a calibrated statement about a bank's loss
distribution. A more restrictive response model, credible experimental
information or justified transport restriction can narrow it. A prediction
for the unobserved band merely names one extrapolation unless its scientific
basis is established.

A discontinuity at the old threshold can, under its design assumptions,
answer a local question about applicants near that threshold. It does not
automatically supply the response throughout a finite newly admitted band.
The assistant must distinguish moving the cutoff by a specified amount
from estimating a local effect at the existing cutoff, and retain the
support or transport argument required to connect them. Public threshold
mechanisms and generated response laws can test this distinction before a
bank considers a costly pilot. The relevant prototype success is correct
translation, appropriate bounds and justified abstention under absent
information, not confident completion of every numerical request.
\label{exec17:l2-policy-end}


% SOURCE docs/release/2026-10-29/source/docs/writing-repair/2026-10-29/src/ef6e4a66d3e6badf.tex CHAR 0
\section[Linguistic measurement and executable interventions]{Linguistic measurement and executable interventions in banking}
\label{sec:exec28-lm}

A language-assisted banking decision begins with an intervention that a bank can actually implement. A complaint team may decide whether to rewrite an acknowledgement, request clarification or escalate a case. An analyst may call the treatment politeness, yet this word can denote the writer's intention, the customer's perception, a classifier score or a particular editing procedure. These objects need not coincide. Their separation is necessary before an agent translates the business question into a graph, proposes an adjustment set or calls an estimator. A well-formed causal query about the wrong object remains the wrong business answer.

\subsection{Intent, perception, labels and actions}

The inspected preprint of \citet{exec28_lm_pryzant} distinguishes writer intent $T$, reader perception $\widetilde T$, a proxy label $\widehat T$ and text $W$. Its first identification result assumes that $W$ blocks the relevant backdoor paths and that intent agrees with perception, $T=\widetilde T$. These are substantive assumptions. A bank may intend reassurance while a distressed customer perceives evasion; a classifier may score deferential wording highly even though the letter makes no useful commitment. Neither a fluent explanation nor high predictive accuracy establishes agreement between these variables.

An operational treatment should identify an editing rule, the eligible message population, the permissible information and the distribution of text versions produced by the rule. For example, a rule might replace a discretionary closing paragraph with a fixed, reviewed clarification of the next servicing step. It must specify whether facts, deadlines and promised remedies can change. If these attributes vary together with tone, the resulting policy effect concerns the entire rewriting procedure. Calling it an isolated effect of politeness requires further restrictions; the proposed product should not silently impose them.

The same distinction guides diagram construction. Pre-existing case severity can influence both message selection and subsequent resolution. A customer response following the edited message is downstream of that action, even if an agent can later extract it into a numerical feature. Putting that response into a baseline adjustment set changes the question and can invalidate identification. The assistant should retain a variable's event time, observation time, decision availability and semantic definition alongside its name. It can suggest candidate arrows and competing explanations, while an explicit causal model states which claims are assumed and what evidence could contradict them.

The source's linguistic formulation provides a conceptual foundation rather than a certificate for an arbitrary embedding. Appendix A discusses the difficulty that a representation can perfectly predict a proxy treatment, leaving no overlap for a within-representation comparison. Its consistency discussion assumes an outcome-sufficient component, appropriate support and consistent regression. The preprint's theorem and surrounding notation alternate perception and proxy symbols; we retain its stated conditions without presenting a silently repaired theorem. A compression method must justify both preservation of adjustment information and usable treatment support. Reducing dimensionality alone proves neither.~\citep{exec28_lm_pryzant}

\subsection{Why classifier accuracy is insufficient}

Under the source's assumptions, define the conditional error probabilities
\begin{align}
 \epsilon_0(W)&=\Pr(\widetilde T=0\mid\widehat T=1,W),\\
 \epsilon_1(W)&=\Pr(\widetilde T=1\mid\widehat T=0,W).
\end{align}
Writing $\delta(W)$ for the conditional perception contrast, its proxy expression has the form
\begin{equation}
 \psi^{\mathrm{proxy}}=E\!\left[\{1-\epsilon_0(W)-\epsilon_1(W)\}\delta(W)\right].
 \label{eq:exec28-lm-proxy}
\end{equation}
The accompanying discussion explicitly qualifies directional attenuation by a better-than-chance classifier and a treatment-effect sign that is homogeneous across possible texts. This qualification must survive translation into a product specification. An overall accuracy number does not ensure that the factor multiplying the contrast is well behaved in each important stratum, or that the underlying effects have a common sign.~\citep{exec28_lm_pryzant}

A separate, fully stipulated measurement example makes the consequence concrete. Let $S\in\{0,1\}$ denote two equally prevalent servicing populations. Within each population, let a latent perception $R$ be assigned with probability one half. Define potential monetary outcomes by
\begin{equation}
 Y(r)=100+\delta_S r,\qquad \delta_0=20,\quad\delta_1=-10,
 \label{eq:exec28-lm-witness}
\end{equation}
in USD per account. Generate an observed proxy $Z$ by independently flipping $R$, conditional on $S$, with probabilities $q_0=9/20$ and $q_1=0$. These are authored assumptions, not estimates from complaints. The proxy is stochastic given the adjustment stratum; this is not the source's deterministic full-text proxy setting or an empirical replication of TextCause.

The true average effect is $(20-10)/2=5$ dollars. Symmetry gives $\Pr(Z=1\mid S)=1/2$. Bayes' rule yields $E[R\mid Z=1,S=s]=1-q_s$ and $E[R\mid Z=0,S=s]=q_s$. Consequently,
\begin{equation}
 E[Y\mid Z=1,S=s]-E[Y\mid Z=0,S=s]=(1-2q_s)\delta_s.
\end{equation}
The two proxy contrasts are $2$ and $-10$ dollars, averaging $-4$ dollars. Overall classification accuracy is
\begin{equation}
 \tfrac12(1-q_0)+\tfrac12(1-q_1)=\frac{31}{40}=77.5\%.
\end{equation}
Exact enumeration of all eight joint cells independently reproduces these conditional means and the algebraic calculation. For a stipulated portfolio of 10,000 accounts, the true gross contrast is USD~50,000 while the proxy comparison suggests minus USD~40,000, before implementation costs.~\citep{exec28_lm_reference}

This example violates the homogeneous-sign condition and therefore does not refute the source's conditional result. A high-accuracy semantic label may still give the wrong direction for a heterogeneous business effect. Evaluation should examine conditional errors, effect heterogeneity, population weighting and support rather than promote a single accuracy score into causal assurance. If perception cannot be measured reliably, the product can report sensitivity ranges or change the intervention to an observable, randomized editing policy. It should not invent a validated perception variable.


\subsection{A useful text signal still needs an intervention argument}
\label{rel66:text-signals}

The distinction between a label and an action also governs claims that a
language model has discovered causes inside case notes. A bank may find that
notes mentioning a payment reminder have higher predicted repayment. Such a
signal can help an analyst search records or formulate a hypothesis. To decide
whether to send a reminder, however, the bank needs the repayment contrast
under an implementable reminder policy and its alternative, using information
available before that choice. A word recorded after the action cannot silently
become a baseline treatment assignment or an admissible adjustment variable.

Three pharmacovigilance studies provide a concrete comparison.
\citet{rel66inferbert} trains a classifier on sentences constructed from
commercially curated FAERS fields. Its two reported tasks concern acute liver
failure and tramadol-related death. Equations 2--6 form groups according to
whether an observed field contains a term, then compare the classifier outputs
in those groups. The paper calls this operation do-calculus. In the inspected
\texttt{causal\_tree} and \texttt{compute\_gap} routines, the predictions are
read from a saved file and partitioned by feature membership; the procedure
does not replace a structural equation or evaluate both actions for the same
unit. Recursive restriction to records containing an earlier term does not
itself supply an adjustment argument. \citep{rel66inferbertcode}

To state what this computes, let $W_i$ denote the recorded text or fields,
$H_i=1$ indicate a selected term's presence, and $f(W_i)$ be a fixed predicted
endpoint probability. On the retained records, the root comparison is
\begin{equation}
 \widehat\Delta_H=
 \frac{\sum_i H_i f(W_i)}{\sum_i H_i}
 -\frac{\sum_i(1-H_i)f(W_i)}{\sum_i(1-H_i)}.
 \label{eq:rel66-score-gap}
\end{equation}
This expression assumes nonempty groups; the code further restricts to
nonmissing fields and groups exceeding its support threshold. Its normal-score
denominator uses the sum of the two within-group sample variances divided by
their respective counts. That calculation describes uncertainty about a
particular score contrast only under its own sampling and fitting assumptions.
It does not establish exchangeability, resolve reporting selection or account
automatically for learning and repeatedly selecting the highest-scoring terms.
The published paper describes adjusted significance levels, while the selected
routine visibly applies an unadjusted one-sided threshold of 0.05. The local
source inspection therefore cannot certify the paper's complete multiplicity
procedure or reconstruct its reported results. \citep{rel66inferbert,rel66inferbertcode}

\citet{rel66deepcausality} extends this design to free text from LiverTox,
using named-entity recognition to propose terms. Its label is inherited from a
drug-level likelihood category: the extracted sentences receive that drug's
positive or negative label. The reported split places 14,361 sentences into
training and test sets in a 9:1 ratio; the account selects its checkpoint using
the test-set classification results. This does not supply an untouched
drug-level evaluation. It leaves both selection optimism and within-drug
dependence to be assessed; no numerical penalty is inferred here. The reported
18 clinically consistent terms out of 20 retained terms measure agreement with
domain knowledge. They do not measure intervention-effect bias or interval
coverage. The paper itself states that the method cannot classify the selected
factors into confounders and colliders. In the pinned notebook, the same score
partition is visible and the significance-filter lines are commented out.
These are reasons to specify the exact implementation and evaluation,
including useful hypothesis-ranking behavior, before adopting its causal
interpretation. \citep{rel66deepcausalitycode}

The subsequent comparison by \citet{rel66pharma} retains both the structured
and free-text tasks. Its domain-expert assessment asks whether enriched terms
are relevant to liver injury. The stated positive predictive quantity is
$\mathrm{TP}/(\mathrm{TP}+\mathrm{FP})$; the narrative reports 0.95 for the
task-adapted DILI-BERT and 0.20 for RoBERTa. Those numbers concern this selected
term-review task. They cannot be read as probabilities that a banking
intervention estimate is correct. The method paragraph cites references 9--10,
whereas the bibliography identifies InferBERT and DeepCausality as 17--18;
the comparison here follows the actual primary methods, rather than treating
the mismatched reference numbers as a new derivation. The source's model-size
and domain-adaptation comparisons remain useful experimental leads within
their reported task and cohort.

An exact authored example isolates the missing step without imitating any of
these experiments. Let unrecorded customer readiness $U$ be Bernoulli with
probability $1/2$, and let a historical action $A$ have probabilities
$\Pr(A=1\mid U=1)=4/5$ and $\Pr(A=1\mid U=0)=1/5$. Suppose repayment $Y$ has
conditional probability $1/10+3U/5$ and is independent of $A$ given $U$.
Independent uniform random variables can implement these two Bernoulli draws,
so intervening on $A$ leaves the repayment equation unchanged. A note records
$A$ exactly. Bayes' rule gives
\begin{equation}
 E[Y\mid A=1]=\frac{29}{50},\qquad
 E[Y\mid A=0]=\frac{11}{50},\qquad
 E[Y(1)-Y(0)]=0.
 \label{eq:rel66-confounding}
\end{equation}
Even the exactly calibrated predictor $f(A)=E[Y\mid A]$ therefore produces
a term-group gap of $9/25=0.36$. For a stipulated USD~100 value per repayment
and USD~10 action cost per account, treating the gap as an effect implies
USD~26 net value; the stipulated intervention law gives minus USD~10. This
is a finite probability construction, not a bank estimate or a replication of
the source studies. Better prediction of this same observational law cannot
choose the causal interpretation.

\subsection{A conditional route from text to adjustment}

Text can nevertheless contain useful adjustment information. The route in
\citet{rel66veitch} begins by assuming that the original text suffices for
the specified causal adjustment. It then asks whether a lower-dimensional
embedding retains the outcome or propensity information needed for estimation.
This addresses loss of already sufficient information during compression; it
does not establish that an arbitrary case note captured all common causes.
Its semi-synthetic experiments test the proposed representations under stated
generators. Its consistency result also requires sufficiently accurate nuisance
models and a well-behaved estimator, with bounded outcomes and overlap given
as sufficient conditions in the discussion. Predictive language performance
alone does not verify these requirements in a bank.

For a single-stage banking use, take $W$ to be text available before a defined
action $A$, let $Z=\lambda(W)$ be its representation, and let $Y$ be the later
outcome. State consistency, conditional exchangeability given $W$, and overlap.
If $Z$ preserves the required adjustment information, define
$m_a(z)=E[Y\mid A=a,Z=z]$ and $e(z)=\Pr(A=1\mid Z=z)$. The target population
then determines the averaging law:
\begin{align}
 \tau_{\mathrm{ATT}}
 &=E[m_1(Z)-m_0(Z)\mid A=1]
   =\frac{E[e(Z)\{m_1(Z)-m_0(Z)\}]}{\Pr(A=1)},\\
 \tau_{\mathrm{ATE}}
 &=E[m_1(Z)-m_0(Z)].
 \label{eq:rel66-target-weights}
\end{align}
The equality for ATT follows from Bayes' rule after identification. It shows
why a representation model and a correct conditional outcome contrast are
insufficient if the final estimator averages over the wrong population.
In the inspected proceedings version, equation 2.1 conditions the outer
expectation on treatment, and equations 2.2--2.3 retain propensity or treated
weights. The later displayed equation 3.1 and the final display in Theorem 3.1
omit those weights or that conditioning while discussing the same target.
We retain the explicit ATT definition and distinguish the displayed formulas;
we do not treat this notation mismatch as a theorem that ATT equals ATE.
The separate appendix proof is not contained in the inspected ten-page
proceedings copy. \citep{rel66veitch}

The distinction can reverse a decision even when all confounding is observed.
Let $Z$ have two equally prevalent strata, with propensities $1/5$ and $4/5$,
and conditional repayment effects zero and $3/10$. One feasible law sets
$m_0(z)=1/5$ in both strata and $m_1(z)=1/5$ and $1/2$, respectively.
Then $\Pr(A=1)=1/2$, the treated population places weights $1/5$ and $4/5$
on the two strata, and
\begin{equation}
 \tau_{\mathrm{ATE}}=\frac{3}{20}=0.15,\qquad
 \tau_{\mathrm{ATT}}=\frac{6}{25}=0.24.
\end{equation}
At USD~100 per repayment and USD~20 per action, the respective net contrasts
are minus USD~5 and plus USD~4 per account in their respective populations.
Neither is a universal policy value: expanding eligibility beyond historical
recipients changes the population and may change support. The example holds
the conditional effects fixed so that the consequence of weighting is visible.

The proposed service should consequently distinguish three outputs: a
term-ranking diagnostic, a measured pretreatment variable, and an identified
effect under a declared adjustment argument. A credible text-assisted study
retains source spans and availability times, evaluates extraction on held-out
customers and periods, and propagates measurement and nuisance-estimation
uncertainty using a sampling design appropriate to repeated notes. It compares
the text representation with structured covariates and simpler text summaries,
reports overlap and sensitivity to omitted information, and uses an untouched
evaluation set after model selection. If the adjustment argument fails,
hypothesis ranking or a randomized prospective action comparison remains a
substantive alternative. These requirements preserve the useful text-processing
contribution without turning its classifier score into a causal assurance.

\subsection{Public complaints as a prototype resource}

The source supplies two distinct kinds of evaluation. Its Amazon setting constructs semi-synthetic outcomes, permitting comparison within the stipulated outcome model. Its CFPB application uses observational complaints and a fast-response indicator defined by a response in fewer than 15 days. It matches complaints with opposite outcomes, selects 4,000 similar pairs and retains the upper and lower quartiles of a politeness score, producing 4,000 selected complaints. These choices define a selected study population, not the entire customer population.~\citep{exec28_lm_pryzant}

The reported effect estimates should not be relabelled as known causal truth. Matching on outcomes and restricting to score extremes require an explicit account of selection and of the population to which a result applies. An agent-generated summary must preserve those choices rather than quote a favourable estimate as a verified benefit of adopting a rewriting service. Generalization to unselected complaints, contemporary servicing and repeated interactions requires additional justification.

For a proposal-stage prototype, public or generated data can establish useful facts at much lower cost than a bank experiment. A frozen public-data snapshot can exercise event-time parsing, missing-text handling, semantic label provenance and the distinction between response time and substantive resolution. A generated outcome construction can test whether an estimator recovers an effect under an explicit assignment and measurement law. An observed complaint archive provides realistic input structure, while a stipulated generator supplies known potential outcomes only within its model. Neither supplies otherwise unobserved counterfactual customer outcomes.

A prospective comparison should freeze the unit, sample selection, source release, permitted fields, intervention rule, estimator and all attempted requests before evaluating the candidate. Customer, case and temporal overlap should be checked when defining training and evaluation partitions. Prompts, retries, manual corrections, failed parses and excluded records contribute to both the denominator and the cost. These are proposed design requirements, not claims that an admitted runtime already satisfies them. A prototype can demonstrate the merit of a later, narrower bank trial without claiming authentic bank validation.

\subsection{From a one-message question to a longitudinal policy}

Repeated servicing introduces feedback. Earlier messages affect later customer responses, which influence subsequent messages and eventual outcomes. The retained CausalCollab work studies sequences of human edits and language-model replies and an incremental stylistic target. Its intervention definition, representations and generated evaluation design should be compared with an executable bank editing policy, rather than treated as interchangeable with a single-stage perception contrast.~\citep{exec13_census_collaboration}

For a separately specified finite-horizon policy, let $H_t$ contain information available immediately before decision $t$, $A_t$ be the editing or escalation action and $g_t(a\mid H_t)$ its prescribed probability. Let $L_{t+1}$ contain the next observations. With a terminal monetary outcome $Y$, the target is $V(g)=E[Y^g]$. Under consistency, sequential exchangeability, positivity for actions used by $g$, and a correctly stated observation process, a finite discrete longitudinal g-formula is
\begin{equation}
 V(g)=\sum_{\bar l,\bar a}m(h_K,a_K)\,p(l_0)
 \prod_{t=0}^{K}g_t(a_t\mid h_t)
 \prod_{t=0}^{K-1}p(l_{t+1}\mid h_t,a_t),
 \label{eq:exec28-lm-policy}
\end{equation}
where $m(h_K,a_K)=E[Y\mid H_K=h_K,A_K=a_K]$ and histories are chronological. Continuous variables require integration. This applies the longitudinal framework; it is not a new identification theorem for arbitrary language.~\citep{hernan2020}

The formula exposes the information the assistant must help specify. The action must denote a procedure, not a classifier label that retrospectively describes a message. Each history must contain the confounders claimed sufficient at that time; unobserved distress or operator judgement cannot be removed merely by generating a diagram. A deterministic action outside historical support invalidates positivity. Representation compression and estimated transition models introduce further approximation and learning questions. The single-stage measurement witness supplies no theorem resolving these longitudinal issues.

The graph-building assistant can produce a reviewable intervention contract: decision times, observable actions, eligible histories, competing graphs, proposed adjustment information, support diagnostics and unresolved assumptions. A reviewer should be able to trace a request to reduce avoidable repeat contact to a defined outcome and a contrast between implementable policies. An unsupported direction, ambiguous treatment or unavailable history should remain visible. This is more informative for prototype investment than a point estimate whose interpretation changed during translation.

\subsection{Adoption choices at the component boundary}

Scientific relevance and executable reuse answer different questions. Three pinned implementations illustrate why component-specific decisions are necessary. The inspected LLM-CD tree contains 2,899 entries and 2,745 files. Its nested BI-CaMML notice states redistribution conditions for that component, while the complete pinned tree has no root notice. Requirements specify Gurobi and pygobnilp versions; compiled artifacts also require correspondence to a reproducible source build. A nested notice and requirements file do not establish permission or a complete installed environment for the whole pipeline. Direct reuse remains deferred pending authoritative component terms, solver entitlement and artifact-to-source provenance.~\citep{exec28_lm_artifacts}

The inspected LLM-and-SCD snapshot has 15 entries and 12 files, without a root licence, README or environment manifest in that tree. Notebook and data paths are discovery leads, but do not establish execution order, full service identities or reproducibility against the paper. The proposal can retain the scientific approach as a comparator while deferring direct reuse until permissions, environment and source-to-result correspondence are available. A separately specified implementation would need its own scope, evaluation and source review.~\citep{exec28_lm_artifacts}

LCM presents a different boundary. Its root code notice is Apache 2.0, while four named pretrained-model links and three experiment-folder links point to external artifacts. Exact model bytes, hashes, provenance and applicable terms have not been established by the inspected code notice. Several dependency versions remain open, and vendored components need their own inventory. The inspected loader sets \texttt{weights\_only=False} and can return \texttt{None} after missing-file or exceptional loads; a caller must distinguish those outcomes from a valid model or a scientifically inconclusive result. An in-tree baseline checkpoint is not a substitute for the advertised LCM weights. These are static source observations, not executed failure rates or legal clearance.~\citep{exec28_lm_artifacts}

For CLI and MCP integration, the proposed request should carry the target and intervention contract, input schema and clock, selected model identity and artifact manifest. The response should separate an unsupported scientific target, an unavailable artifact, an execution failure, an inconclusive analysis and a successful result with declared assumptions. A fallback may be proposed as an explicit new method choice; silently replacing the intended model changes the comparison. A public fixture should test these branches and record all attempted requests and costs. Where external artifacts remain unavailable, independent work on an authorized estimator, public-data design or semantic contract can proceed. The proposed product must preserve the distinction between a meaningful scientific question, a justified answer and an executable method.


% SOURCE docs/release/2026-10-29/source/docs/writing-repair/2026-10-29/src/5cbfee0b972f4c9c.tex CHAR 0
\section[From business phrases to causal variables]{From business phrases to causal variables and defensible diagrams}
\label{sec:exec28-cm}

A useful causal assistant must decide what the variables mean before it can reason about arrows. In a servicing discussion, current borrower, delinquent borrower and financially distressed borrower may denote values of different operational definitions, overlapping classifications or measurements taken at different times. Automatically drawing one node for each phrase can turn a vocabulary mistake into a scientific model. The first task is to propose a coherent representation and preserve alternatives that cannot yet be resolved.

\subsection{Two concrete language tasks}

The inspected preprint of \citet{exec28_cm_razouk} evaluates two tasks directly relevant to this problem. The first asks whether two textual entities represent different values of the same variable, without supplying that variable's name. The second asks whether an entity refers simultaneously to values of several supplied variables. The source calls the second object an interaction entity. These are useful semantic tasks, distinct from verifying a causal arrow, identifying an effect or establishing an intervention's feasibility.

For a bank, consider whether payment within the grace period and payment after the grace period are values of a common timing variable. The answer depends on the account, payment obligation, time origin and grace-period definition. A phrase describing a mortgage customer with a payment holiday combines product and servicing status. Detecting that combination does not establish that these statuses interact in their causal effect on default. The assistant should first return candidate variable definitions, value domains, unit identifiers, observation clocks and textual evidence. A subsequent model specifies relationships among those variables.

A proposed representation can be written as
\begin{equation}
 \mathcal V_j=(\text{unit},\text{quantity},\text{domain},\text{time},
                \text{measurement},\text{availability}),
\end{equation}
where a phrase may map to a value, a predicate on a domain or an unresolved collection of variables. This is a product design convention, not a theorem that language has a unique causal decomposition. A high-risk label can denote a score category, a decision-maker's belief or an unobserved propensity. Treating them as interchangeable can make a score's prediction appear to be an intervention on underlying risk.

Pairwise language answers also need consistency review. Suppose an assistant groups phrase $e_0$ with $e_1$, groups $e_1$ with $e_2$, and rejects grouping $e_0$ with $e_2$. Those answers cannot all define an equivalence relation for a single partition. The proposed product should expose the conflicting triangle and relevant definitions. Taking the transitive closure silently can merge incompatible variables, while deleting one pair without a reason hides the conflict. Multiple candidate representations are preferable when the language remains ambiguous. The independent finite reference reproduces the contradiction.~\citep{exec28_cm_reference}

\subsection{What the evaluation actually establishes}

The source uses a within-subjects design in which seven named language models generate examples and evaluate examples from themselves and other models. It sets temperature to zero and uses the same English prompts. Its domain tables contain eight domains, including finance. Task 1 uses 560 positive and 560 negative examples; Task 2 uses 559 positive and 560 negative examples. These are reported evaluation denominators, not independently audited bank expert labels.~\citep{exec28_cm_razouk}

The results illustrate why a single model ranking is insufficient. In the source's Task 1 table, GPT-4-turbo has reported precision 97\%, recall 41\% and kappa 40\%; rankings differ for the second task. These historical labels and reported numbers describe that generated-data evaluation. They are not a recommendation for a current provider, a service guarantee or evidence that 97\% of suggested banking variables are correct. High precision with low recall can leave important distinctions unresolved. The cost of an omitted variable may differ greatly from that of a false proposed merge.~\citep{exec28_cm_razouk}

The paper explicitly states that its numbers derive from generated data and need comparison with domain-specific expertise. Agreement with generation-derived labels can reward shared conventions or mistakes. Cross-model evaluation broadens the comparison but does not produce independently established causal truth. Temperature zero supplies neither calibrated uncertainty nor an immutable service implementation. Domain dependence is a reason to stratify later assessment by banking journey, variable type, language and ambiguity rather than transport an overall score to every workflow.

A proposal-stage evaluation can retain inexpensive generated fixtures alongside a small, independently adjudicated public set. Generated fixtures make intended variable/value relationships explicit and support controlled perturbations. Public documents provide realistic ambiguity, missing definitions and conflicting timestamps. Their reference labels need documented provenance; disagreement should be preserved instead of resolved by asking the same model to grade itself. The output includes an unresolved option, with performance reported separately for accepted, rejected and deferred cases. Coverage, incorrect merges and splits, unsupported certainty, human review effort and full workflow cost all matter.

\subsection{A compound phrase does not establish causal interaction}

Let $A$ denote offering a payment holiday and $B$ an outreach action. A phrase mentioning both actions does not determine their additive interaction in a monetary outcome. That requires a four-regime contrast,
\begin{equation}
 I=E[Y(1,1)]-E[Y(1,0)]-E[Y(0,1)]+E[Y(0,0)].
 \label{eq:exec28-cm-interaction}
\end{equation}
In an authored example, set $Y(a,b)=100+10a+20b$ dollars per account. The outcomes are 100, 110, 120 and 130, giving $I=0$ despite the compound phrase. Changing the stipulated law to $100+10a+20b+5ab$ changes the joint cell to 135 and gives $I=5$ dollars. Exact enumeration checks both contrasts. This is a scale-specific potential-outcome calculation, not measured servicing interaction.~\citep{exec28_cm_reference}

A deterministic compound node $C=AB$ may be convenient, but the analyst must state what an intervention on $C$ means. Setting $C=1$ without specifying $A$ and $B$, or changing $C$ while holding its defining parents fixed, need not describe a feasible action. Multiple versions of a compound label can have different consequences. The assistant should require an intervention definition rather than infer it from a noun phrase.

The source invokes causal transitivity as motivation. That terminology should not become a numerical rule that every directed path has a nonzero total effect. Let $U$ be uniform on $\{0,1,2,3\}$, define $B=\mathbf 1\{U<1+2A\}$ and set $Y=100+20B-10A$. The model has an $A\to B\to Y$ path and a direct $A\to Y$ path. Under interventions on $A$, the mediator probability changes from $1/4$ to $3/4$. This produces a mediated mean change of USD~10 and a direct change of minus USD~10. Both intervention means are USD~105: the total mean effect is zero through cancellation, although both mechanisms are present.~\citep{exec28_cm_reference}

This finite example distinguishes graph structure, mechanisms and a numerical contrast; it does not assert that the source's experiment is invalid. The proposed validator should check stated graph semantics, temporal order and intervention definitions, rather than impose a blanket effect-magnitude rule. Statistical dependence also cannot by itself endorse a causal arrow. A graph, an observation law and an estimand serve different purposes.

\subsection{An integration contract for semantic proposals}

The proposed CLI or MCP interface can expose variable construction as a reviewable artifact. Each variable should carry its source span, value domain, population, event clock, availability clock, measurement rule and possible action semantics. A mapping can be accepted, disputed or unresolved. Relations should distinguish source assertions, model-generated hypotheses, statistical associations and approved modeling assumptions. A syntactically valid DAG must not erase these distinctions.

The retained source links to an anonymous code-and-data repository. The link is an acquisition lead; it does not establish a pinned release, licence, model service identity or executable package in this workspace. Direct runtime adoption is deferred. The tasks remain useful leads for an independently specified prototype. Original prompts and reference labels are separate artifacts whose rights and provenance require review. A prototype using separately authorized fixtures can test representation quality without claiming to reproduce the source's runtime.

Before estimation, the assistant should present competing variable definitions, an intervention and outcome, unresolved timing or measurement issues, and the assumptions needed to compare selected policies. The business owner can then see what information is missing and why it matters. The result is a traceable scientific specification whose ambiguity remains visible, rather than a diagram that merely looks complete.


% SOURCE docs/release/2026-10-29/source/docs/writing-repair/2026-10-29/src/fd53384e3b953ff7.tex CHAR 0
\section[Causal tasks and representation interventions]{Causal tasks and representation interventions: what the evidence actually supports}
\label{sec:exec28-fr}
A causal assistant must preserve the distinction between answering a causal question and changing a language model.
The distinction is practical for a bank. An assistant might extract covariates from a scanned credit file,
propose a graph for a servicing decision, identify an admissible adjustment set, or estimate the effect of a payment reminder.
Better image descriptions can improve the first step without establishing the assumptions needed by the last three.
Conversely, a correct identification formula cannot repair an incorrectly extracted treatment date. Each part needs its own
target, observation contract, evidence and failure state.

\subsection{A task taxonomy that retains the information supplied}
Ma's selected September 2024 survey separates pairwise and full-graph discovery, effect estimation,
attribution, counterfactual reasoning and explanation. It discusses prompting, fine-tuning, combinations with conventional
causal algorithms, and knowledge or tool augmentation. This is useful as a taxonomy, rather than a new common benchmark.
The survey explicitly notes that the evaluations collected in its comparison tables can use different prompts and settings;
its dataset table also distinguishes pairs, relations, samples and questions as units. A percentage attached to one of these
units must not silently become a percentage of successfully completed banking analyses.~\citep{exec28_fr_ma}

Consider a hypothetical collections analyst asking whether a reminder reduces thirty-day arrears. A pairwise query about
reminders and arrears may recover familiar domain associations. A graph task also needs eligibility, prior delinquency,
contactability, exposure opportunities and the decision clock. An effect task needs an intervention policy and a population;
a counterfactual question about one borrower requires more than the marginal average effect. These are different scientific
requests even if the interface presents each as a sentence. The proposed compiler should retain the supplied graph,
assumptions, empirical information and requested quantity as separate fields. A model that receives the correct graph and
potential-outcome table has solved a different problem from one asked to infer both from administrative records.

The survey's discussion identifies prompt sensitivity, memorized knowledge, inconsistent answers and scenario-dependent
rankings; it does not establish a universally best model. Its discussion of high-dimensional text also flags support problems.
For a longitudinal bank example, encoding all earlier messages does not by itself create treatment overlap at each history,
nor establish that every common cause was recorded before the action. Such omissions should produce a restricted answer,
an identified set or a refusal under the library's existing contracts. Additional prose from an agent does not supply the
missing observation law.~\citep{exec28_fr_ma}

For public prototype evaluation, separate three modes: supplied-structure symbolic queries with known answers;
discovery under a declared generative model; and interpretation of genuinely observed documents with independently checked
labels. Keep the data-generating mechanism and test answers outside the evaluated agent's context. Record rejected,
unidentified, timed-out and malformed requests in the denominator. A public corpus can test extraction and workflow usability,
while generated outcomes test a stipulated causal target. Neither alone validates a bank's assignment mechanism.
This is a proposed evaluation design, not a claim that all these tests have been executed.

\subsection{Causality used inside a multimodal model}
Hu and colleagues' Causal-LLaVA preprint addresses object hallucination. It adds a projector module and a module in the
language model's final transformer block, using visual and textual dictionaries formed from category-averaged representations
of 5,000 instruction samples. Their proposed graph and approximate adjustment motivate the architecture; the reported
experiments evaluate image-language behavior. They do not evaluate identification of a servicing or lending effect.
The inspected version reports training on eight H20 GPUs and an additional dictionary-estimation cost of approximately
one hour. These are the authors' experimental conditions, not a measured cost for this proposal.~\citep{exec28_fr_hu}

The source's tables support a qualified empirical comparison. Its LLaMA-2-7B comparison reports CHAIR sentence hallucination
33.0 versus 30.9 and recall 65.6 versus 65.9; these are distinct metrics. It also reports small regressions in particular
general-comprehension entries: the 13B ScienceQA entry changes from 65.1 to 65.0, and the Vicuna-v1.5-7B VizWiz entry
from 50.1 to 49.5. The blanket reading that every score improves is therefore unwarranted. The table does not by itself
establish significance, cross-bank transfer or an end-to-end error rate. The source's limitations mention computing demands
and confounder sensitivity across environments and datasets.~\citep{exec28_fr_hu}

\subsection{Three mathematical distinctions behind the architecture}
Under an appropriate adjustment set and the usual identification conditions, the target is an arithmetic mixture
\[
 p(y\mid\operatorname{do}(x))=\sum_z p(y\mid x,z)p(z).
\]
Its weighting distribution is the target marginal, not automatically uniform. A uniform weighting is an additional target
choice. Nor does the formula license replacing the marginal with input-dependent attention weights without proving that the
result retains the same intervention meaning. The Causal-LLaVA appendix explicitly conditions dictionary weights on hidden
states to preserve representational capacity. That may be a useful predictive architecture, but predictive usefulness and
identification of the displayed intervention target require separate arguments.~\citep{exec28_fr_hu,pearl2009}

The normalized weighted geometric mean is likewise distinct from the arithmetic mixture. Let the two conditional event
probabilities be \(1/10\) and \(9/10\), with target weights \(3/4,1/4\). The arithmetic probability is \(3/10\).
For a binary softmax representation, the corresponding logits are \(-\log9,\log9\); their weighted mean is
\(-\tfrac12\log9=-\log3\), whose logistic probability is \(1/4\). Thus even a linear logit map does not make the arithmetic
mixture equal to the normalized geometric approximation. The difference is \(1/20\), or five percentage points.
This finite construction tests a mathematical equality, not the authors' observed benchmark performance.~\citep{exec28_fr_reference}

The same stipulated event law illustrates the importance of weights. Uniform weights give \(1/2\); input-conditioned
weight \(4/5\) on the high-risk category gives \(37/50\). If the event incurs a stipulated USD 1,000 loss, the four
calculations produce expected losses of USD 300, 250, 500 and 740, respectively. These amounts are exact consequences of
the example's assumptions. They are not bank forecasts, a superiority result, or evidence that a particular implementation
actually uses these example weights.~\citep{exec28_fr_reference}

A second distinction concerns joint confounders. In general, marginalizing visual and textual confounders requires their
joint law, conditional or interventional as appropriate. Replacing that law by a product of marginals requires an independence
condition. If two Bernoulli variables are identical and each is one with probability \(1/2\), their joint-one probability is
\(1/2\), whereas multiplying the marginals gives \(1/4\). Equation (4) in the inspected source uses a product of conditional
marginals. The proposed product should require the corresponding independence premise or preserve the joint distribution;
the example establishes the consequence of omitting that premise, not a measured failure of the trained network.
Freezing an encoder also does not by itself prove that contextual common causes disappear from its representations.
That causal claim needs a model of the data and the frozen transformation.~\citep{exec28_fr_hu,exec28_fr_reference}

\subsection{A bounded adoption decision and interface consequence}
The appropriate proposal disposition is to retain the architecture as a candidate measurement component and defer direct
adoption pending exact artifact, weight, data and dependency rights and a reproducible matched evaluation. The source provides
a repository lead; this review has not pinned or executed that repository. A low-cost first prototype should compare against
the existing extractor on a public or generated document set with separately adjudicated object labels, then assess how
extraction errors propagate into a fixed causal query. A successful vision score is neither an effect-identification certificate
nor a substitute for an uncertainty calculation.

CLI and MCP requests should therefore distinguish measurement improvement, model-mechanism intervention and a real-world
policy intervention. Return the target population, dictionary construction, weighting rule, approximation status, graph
assumptions and unavailable artifacts explicitly. Report a source-derived method as source-derived, a proposed adapter as
proposed, and an executed test with its exact code and data identities. For the banking analyst, the useful output is an
auditable answer whose scientific target survives every translation and approximation.


% SOURCE docs/release/2026-10-29/source/docs/writing-repair/2026-10-29/src/2f8f1fc13e8d1b65.tex CHAR 0

\section{Clinical graph refinement as a test of business-to-causal translation}
\label{sec:exec25-ecmo}
A bank translating a servicing problem into a causal diagram must distinguish
an informed graph proposal from an identified intervention. The clinical study
of Mathur and collaborators is useful because it combines approximate language-model
knowledge, expert constraints and a small observational dataset in a concrete
workflow. Its value for this proposal lies in the questions it exposes: what is
the time of each observation, which edges are forbidden by substantive knowledge,
which graph differences matter for the decision, and what constitutes independent
evidence that refinement helped? The inspected source is the five-page author
original in a dated repository, not an asserted byte-equivalent copy of the
six-page publisher layout. The empirical results below are the authors' reports;
the banking examples and acceptance requirements are project constructions.
\citep{exec24_ecmo_original,exec24_ecmo_reference}

\subsection{What the hybrid procedure contributes}
The study uses 71 patients from one center and eight Boolean variables. Seven
risk indicators summarize the first 24 hours of extracorporeal membrane
oxygenation, while neurological injury is determined using tests at the end of
the run. Experts select the variables and forbidden directions; four named
language-model families supply approximate initial graph knowledge. The workflow
first deletes edges and then permits full refinement through local score search.
It evaluates proposed graphs against an expert graph using spurious directed
edges, structural Hamming distance (SHD) and structural intervention distance
(SID). These are comparisons with the selected reference graph. They do not
turn the reference into observed intervention truth. The paper itself notes
the small single-center sample, possible lack of causal sufficiency and the need
to extend the method across time.~\citep{exec24_ecmo_original}

For a bank, the corresponding task is to translate a question such as whether
an early hardship offer reduces later delinquency into a variable dictionary,
an intervention, a population and a measurement clock. A proposed node must say
whether it is a pre-offer condition, the offer, uptake, a later response or a
selection event. The expert's prohibition against a later event causing an
earlier measurement is useful information. It cannot, by itself, resolve
unmeasured common causes, selection into the dataset or the direction between
two events collapsed into the same interval. We therefore retain the hybrid
procedure as a reference for elicitation and constrained hypothesis generation,
with a separate identification decision for each proposed banking query.
\citep{exec24_ecmo_reference}

The expert restrictions define an admissible diagram set $\mathcal{G}(K)$.
Here $K$ records both the timed restrictions and their provenance. A proposed
local refinement has the form
\begin{center}
$\widehat G=\operatorname{LocalSearch}_{G\in\mathcal{G}(K)}
S(D,G;G_0)$,
\end{center}
with data $D$, initial graph $G_0$ and score $S$. This notation
describes a search output, not a proof that $\widehat G$ identifies the effect
of an offer. A public test must hold out both cases and the expert material used
for evaluation from the elicitation step. Otherwise agreement can reflect
reused constraints. If several admissible graphs imply different answers, the
interface should retain the ambiguity or refuse the requested point estimate.
The scientific oracle must judge that behavior separately from whether the
agent returns a syntactically valid graph.~\citep{exec24_ecmo_reference}

\subsection{Refinement does not uniformly improve the reported comparisons}
Table 1 of the author original reports the following union-graph means. Lower
values represent fewer disagreements under the stated reference comparison.
\citep{exec24_ecmo_original}
\begin{center}
\begin{tabular}{lrrr}
\toprule
Stage & Spurious edges & SHD & SID \\
\midrule
Initial union & 9 & 9 & 5 \\
Deletion only & 3 & 5.4 & 5.1 \\
Full refinement & 4.8 & 6.6 & 6.2 \\
\bottomrule
\end{tabular}
\end{center}
Thus full refinement raises the three reported means relative to deletion only
by $1.8$, $1.2$ and $1.1$, respectively. Deletion reduces the union's edge and
SHD comparisons but does not reduce its SID mean from the initial value.
For LLaMA in Table 2, SHD is $5$ initially, $6.1$ after deletion and $8.1$
after full refinement. The authors discuss plausibility of some added edges
following expert inspection. This is a distinct observation from improved
reference-graph scores. These descriptive differences neither establish a
paired significance result nor show that every refined edge is causally wrong.
They do rule out presenting the reported table as uniform metric improvement.
\citep{exec24_ecmo_original,exec24_ecmo_reference}

The product comparison should therefore expose initial, deletion-only and full
outputs with the same cases, constraints and scorer. The banking choice among
them should depend on a preregistered intervention query and asymmetric costs
of an unsupported edge, an omitted edge and an appropriate refusal. A method
that improves one structural distance can still change a decision-relevant
adjustment set adversely. A better in-sample score is evidence about the search
criterion; causal utility requires the separately defined oracle. These are
requirements for the proposed public prototype, not reported clinical or bank
experimental findings.~\citep{exec24_ecmo_reference}

\subsection{A measurement clock that the diagram cannot recover}
Consider two events $A$ and $B$ in the first day of a servicing case. Define
$A^*=\mathbf{1}\{A\text{ occurs within 24 hours}\}$ and define $B^*$
similarly. One constructed history has $A$ at hour 2 and $B$ at hour 20; a
second reverses those times. Both map to $(A^*,B^*)=(1,1)$. The map from
histories to summaries is many-to-one, so these indicators alone cannot
recover the within-window ordering. This example establishes an information
loss, not an impossibility theorem for every causal effect that might be
identified with further assumptions or additional information.
\citep{exec24_ecmo_reference}

In a bank, an indicator for both a missed payment and a hardship contact during
one month could conceal whether the contact preceded the missed payment. The
translator should ask for event time, observation time, action eligibility and
the decision deadline before proposing a causal ordering. If only summaries
exist, it should record which directions depend on external knowledge and
carry that uncertainty into the query. A later outcome cannot be promoted
to a pre-action adjustment variable merely because all columns occupy one
table. The clinical paper's early-window construction motivates this design
question; it does not validate a specific bank's event chronology.
\citep{exec24_ecmo_original,exec24_ecmo_reference}

\subsection{Source semantics and complete experimental denominators}
The pinned implementation provides concrete inspection targets. The selected
commit contains deletion-only and full search routes, acyclicity and forbidden
edge checks, and a structured score that compares a tabular score with fitted
tree alternatives using the supplied data. The inspected tree-fitting path
does not supply a held-out validation claim. The metric routine computes
directed spurious edges by set difference and calls SHD and SID against the
expert graph; selected score variants and graph conversions must remain part
of a reproducible profile. Root GPLv3 text and declared requirements were
inspected, but this does not clear the rights or compatibility of an entire
installed distribution. No acquired code was executed.
\citep{exec24_ecmo_code}

The metrics routine initializes ten slots to zero, fills slots from filenames
and later takes their mean and population standard deviation. The inspected
path contains no seen-file completeness test. A conditional arithmetic example
shows why a trial harness needs one: if six files each have score 8 while four
slots retain zero, averaging ten slots gives $48/10=4.8$, while averaging the
six observed scores gives $48/6=8$. Neither rule should be silently substituted
for a complete trial. The example does not assert that any actual result file
was missing: the archived result contents were not inspected. A prototype
must distinguish absent, failed, refused and successfully scored cases, report
all their counts, and prevent a default zero from being interpreted as perfect
agreement.~\citep{exec24_ecmo_code,exec24_ecmo_reference}

The filename parser also takes one character from the sample identifier.
Constructed names containing sample IDs 1 and 10 collide under that rule.
The source explicitly limits its supported configuration to ten bootstrap
samples; the collision example is a warning against extending that profile
without changing and testing identity handling, not an observed failure within
the stated configuration. A reusable worker should use explicit case and
replicate IDs, reject duplicates, validate the expected inventory and test
edge orientation, variable ordering and empty outputs. Native conversion of
FCI endpoint representations and the installed metric dependencies remain
separate test obligations.~\citep{exec24_ecmo_code,exec24_ecmo_reference}

\subsection{Public qualification and a bounded reuse decision}
The next defensible experiment uses public or generated cases with known
structural mechanisms, timed observations and held-out intervention answers.
The case inventory should vary measured and hidden common causes, selection,
measurement aggregation and supported graph classes. Freeze the model and
prompt identities, initial constraints, search settings, metric definitions
and train/test split before comparing expert-only, data-only, LLM-only,
deletion-only and full-refinement routes. Report graph agreement, identified
query correctness, calibrated uncertainty where justified, refusal behavior
and complete execution costs separately. A public test can assess these
proposed behaviors without access to the omitted patient dataset. It cannot
reproduce the clinical findings or establish performance on authentic bank
assignment mechanisms.~\citep{exec24_ecmo_reference}

For budgeting, an explicitly illustrative total of USD 1,800 across 100
submitted requests is USD 18 per submission. If only 60 produce scientifically
usable results, the same total is USD 30 per usable result. The numerator
must include expert work, unsuccessful attempts, model calls, infrastructure,
licensing and adjudication attributable to that experiment. These are exact
cost-denominator witnesses, not observed prices or an investment return.
Value requires a separately specified decision loss and matched comparator.
\citep{exec24_ecmo_reference}

Our present decision is to use the inspected work as a methodological reference
for constrained graph elicitation and as a source of explicit prototype tests.
Direct adoption or a whole-library port awaits a supported public profile,
dependency and rights review, conversion tests and independent scientific
evaluation. The final publisher layout, supplementary prompts and model
versions, and archived result contents remain uninspected. The author
original, pinned source and exact project examples support the limited
comparison above; they do not establish an implemented or accepted banking
capability.~\citep{exec24_ecmo_original,exec24_ecmo_code,exec24_ecmo_reference}

\subsection{What the four translations require of their measurements}

The applications show why translation must reach the decision consequence.
Additional deposits can increase funding cost after repricing; a first action
has different value under different later policies; and a favorable mean
among approved applicants does not determine the value of expanding approval.
Recorded notes and extracted labels introduce separate measurement and
selection questions. A graph or a fluent description cannot restore the
information lost by an inappropriate aggregation or target substitution.

The component comparisons make the same requirement operational. Perceived
politeness, an editing rule and a classifier label are different objects.
Compound phrases do not establish interaction effects, and improved graph or
image scores do not establish an intervention value. The next chapter turns
to the economic variables themselves: how to define their units, measurements,
relationships and clocks so that the formal question remains the one the bank
intends to answer.

